\documentclass[11pt,a4paper,twoside]{tesis}
% SI NO PENSAS IMPRIMIRLO EN FORMATO LIBRO PODES USAR
%\documentclass[11pt,a4paper]{tesis}

\usepackage{graphicx}
\usepackage[authoryear]{natbib}
\usepackage{adjustbox}
\usepackage{amssymb}
\usepackage{amsmath}
\usepackage{amsthm}
\usepackage{booktabs}
\usepackage[tables]{xcolor}
\usepackage{easy-todo}
\usepackage{colortbl}
\usepackage[utf8]{inputenc}
\usepackage[spanish]{babel}
\usepackage[left=3cm,right=3cm,bottom=3.5cm,top=3.5cm]{geometry}

\begin{document}

%%%% CARATULA
% Comentar y descomentar según corresponda
\def\titulo{Licenciado }

\def\autor{Juan Manuel Pérez}
\def\tituloTesis{Mimetización entre interlocutores}
\def\runtitulo{Medición de la mimetización entre interlocutores utilizando series de tiempo}
\def\runtitle{Measuring entrainment between speakers using time series}
\def\director{Agustín Gravano}
\def\codirector{Ramiro Gálvez}
\def\lugar{Buenos Aires, 2015}
\input{caratula}

%%%% ABSTRACTS, AGRADECIMIENTOS Y DEDICATORIA
\frontmatter
\pagestyle{empty}
\chapter*{\runtitulo}

Los sistemas de diálogo humano-computadora son cada vez más frecuentes, y sus aplicaciones comprenden una amplia gama de rubros: desde aplicaciones móviles, motores de búsqueda, juegos, hasta tecnologías de asistencia para ancianos y discapacitados. Si bien es cierto que estos sistemas logran captar buena parte de la dimensión lingüística de la comunicación humana, tienen un déficit importante a la hora de procesar y transmitir el aspecto superestructural de la comunicación oral, que radica en el intercambio de afecto, emociones, actitudes y otras intenciones de los participantes.

El \emph{entrainment} (mimetización) es un fenómeno inconsciente que se manifiesta a través de la adaptación de posturas, forma de hablar, gestos faciales y otros comportamientos entre dos o más interactores. A su vez, la ocurrencia de esta mimetización está fuertemente emparentada con el sentimiento de empatía y compenetración entre los participantes. En nuestro caso, nos es de interés el \emph{entrainment} sobre las variables acústico-prosódicas, como el tono, intensidad, y otras.

En el presente trabajo, nos proponemos explorar y refinar una métrica del entrainment acústico-prosódico definida en trabajos previos. Analizamos la relación entre los valores obtenidos y las percepciones sociales que terceros tienen sobre las conversaciones, en un corpus de diálogos orientados a tareas en inglés.

\bigskip

\noindent\textbf{Palabras claves:} Procesamiento del Habla, Series de Tiempo, Entrainment, Regresión Lineal

\input{abs_en.tex}

\tableofcontents

\mainmatter
\pagestyle{headings}

%%% Macros
\newcommand{\absentrainment} {\emph{unsigned entrainment}\xspace}
\newcommand{\fwentrainment}[1] {\mathcal{E}_{#1}}

\newcommand{\entrainment} {\emph{entrainment}\xspace}
\newcommand{\disentrainment} {\emph{disentrainment}\xspace}
\newcommand{\ap} {acústico-prosódicas\xspace}

\newcommand{\TAMA} {\emph{TAMA}\xspace}
%%% Variables A/P
\newcommand{\apvar}[1]{\emph{#1}\xspace}

\newcommand{\ENGMAX} {\apvar{Int Max}}
\newcommand{\ENGMEAN} {\apvar{Int Mean}}
\newcommand{\FOMAX} {\apvar{F0 Max}}
\newcommand{\FOMEAN} {\apvar{F0 Mean}}
\newcommand{\NOISETOHARMONICS} {\apvar{NHR}}
\newcommand{\SYLAVG} {\apvar{Sílabas/seg}}
\newcommand{\PHONAVG} {\apvar{Fonemas/seg}}
\newcommand{\LOCALSHIMMER} {\apvar{Shimmer}}
\newcommand{\LOCALJITTER} {\apvar{Jitter}}


%%% variables sociales

\newcommand{\socialvariable}[1] {\emph{#1}\xspace}
\newcommand{\svcontributes} {\socialvariable{contributes-to-completion}}
\newcommand{\svclear} {\socialvariable{making-self-clear}}
\newcommand{\svengaged} {\socialvariable{engaged-with-game}}
\newcommand{\svplanning} {\socialvariable{planning-what-to-say}}
\newcommand{\svencourages} {\socialvariable{gives-encouragement}}
\newcommand{\svdifficult} {\socialvariable{difficult-for-partner-to-speak}}
\newcommand{\svbored} {\socialvariable{bored-with-game}}
\newcommand{\svdislikes} {\socialvariable{dislikes-partner}}

%%% Time series

\newcommand{\varnorm}[1] {
    (#1_t - \mu_{#1})
}

%%% Regresión and stuff
\newcommand{\estslope} {\widehat{\beta_2}}
\newcommand{\estintercept} {\widehat{\beta_1}}
\newcommand{\softhl} {\rowcolor[gray]{.88}}
\newcommand{\hl} {\rowcolor[gray]{.80}}
\newcommand{\stronghl} {\rowcolor[gray]{.75}}

% psl is "Positive Slope"
\newcommand{\psl} { $+$ }
\newcommand{\ppsl} { $++$ }
\newcommand{\pppsl} { $+++$ }

% nsl stands for "Negative SLope"
\newcommand{\nsl} { $-$ }
\newcommand{\nnsl} { $--$ }
\newcommand{\nnnsl} { $---$ }


%%%% ACA VA EL CONTENIDO DE LA TESIS
\chapter{Introducción}
\begin{frame}
  \frametitle{Sistemas de diálogo Humano-Computadora}

  \begin{enumerate}[<+->]
    \item Sistemas actuales
    \item Bien en la parte lingüística de la comunicación: entender y transmitir mensajes estructuralmente correctos.
    \item Mal en la parte superestructural: intercambio de emociones, actitudes, etc.
    \item El presente trabajo trata de hacer un (pequeño) aporte sobre el análisis de la ``naturalidad'' de las conversaciones.
  \end{enumerate}
\end{frame}


\begin{frame}
  \begin{columns}
    \column{0.5\textwidth}
    \frametitle{Ejemplos de ``falta de naturalidad''}
    \begin{enumerate}
      \item Sistemas de llamadas comerciales
      \item Siri, Google Now
      \item Otros?
    \end{enumerate}
    \column{0.5\textwidth}
    \begin{figure}
      \includegraphics[scale=0.25]{images/hal.jpg}
    \end{figure}
  \end{columns}
\end{frame}


\begin{frame}
  \frametitle{Prosodia}
  \framesubtitle{Definir acá qué es la prosodia aproximadamente}

\end{frame}



\begin{frame}
  \frametitle{Entrainment}

  \begin{enumerate}
    \item Fenómeno ubícuo en la comunicación
    \item Ocurre a varios niveles
    \item Largamente estudiado en psicología de comportamiento (referencias)
    \item (ya está, la próxima conversación que tengan afuera de acá van a chequear ésto)
  \end{enumerate}
\end{frame}

\begin{frame}
  \frametitle{¿Y cómo lo medimos?}
  La definición de entrainment hasta acá vista es heurística! ¿Cómo definimos una medida para esto?

  Vamos a explorar una métrica definida en trabajos anteriores, pulirla un poco, y verificar que efectivamente capture ciertas características del entrainment.

  ¿Cómo? Usando un corpus con anotaciones sociales


\end{frame}



\chapter{Antecedentes}
\begin{frame}
  \frametitle{Otras métricas}

  \begin{itemize}
    \item La mimetización es un fenómeno tanto lineal: se va acentuando a lo largo de la conversación
    \item Pero también es un fenómeno dinámico: va variando localmente a lo largo de la conversación.
  \end{itemize}

  Muchas métricas sólo toman la parte global, dividiendo la conversación en 2 o más partes y luego calculando la diferencia entre las medias de las diferentes variables acústicas en cada sección.
\end{frame}


\begin{frame}
  \frametitle{Problema del alineamiento de tiempo}

  \begin{figure}[t]
    \includegraphics[scale=0.40]{images/conversation_turns.pdf}
  \end{figure}

  \begin{itemize}
    \item ¿Cómo comparamos los diferentes turnos de una conversación?
    \item Comparar uno a uno es un enfoque simplista y no representativo de la realidad
  \end{itemize}
\end{frame}

\begin{frame}
  \frametitle{Método TAMA}
\end{frame}




\chapter{Método y materiales}
\section{Columbia Games Corpus}

\newcommand{\objectgame} {\emph{Juego de objetos}}


Empleamos el Columbia Games Corpus  \cite{GRAV2009} que consiste de doce conversaciones diádicas (i.e., con dos participantes) entre trece personas angloparlantes distintas. Todos los participantes reportaron hablar inglés americano estándar, y no tener problemas de audición. La edad de los participantes se encuentra en el rango de los 20 a 50 años.

Las grabaciones se hicieron en 44 kHz, 16 bits con un canal separado para cada hablante; luego fueron guardadas en 16 kHz para el presente estudio. Cada sesión duró aproximadamente 45 minutos, totalizando 9 horas de
diálogos, 70.259 palabras (2.037 únicas) para todo el cuerpo de datos. Todas las conversaciones cuentan con transcripciones textuales alineadas temporalmente a la señal de audio, realizadas por personal especialmente entrenado.

En cada sesión, se sentó a dos participantes (quienes no se conocían previamente) en una cabina profesional de grabación, cara a cara a ambos lados de una mesa, y con una cortina opaca colgando entre ellos para evitar la comunicación visual. Los participantes contaron con sendas computadoras portátiles conectadas entre sí, en las cuales jugaron una serie de juegos simples que requerían de comunicación verbal. El primero de ellos es un juego de cartas que no consideramos en el presente estudio por tratarse esencialmente de monólogos o diálogos con poca interacción. Luego de esto, pasaron al juego que analizamos, denominado `juego de objetos'.

\subsection{Juego de Objetos}

En el juego de objetos, la pantalla de cada jugador mostró un tablero con varios objetos, entre 5 y 7, como se ve en la Figura \ref{objects_game}.
Para uno de los jugadores (el Descriptor) el objeto \emph{Objetivo} aparecía en una posición aleatoria entre otros objetos. Para el otro jugador, a quien llamaremos el Seguidor, el objetivo aparecía en la parte baja de la pantalla. Entonces, al Descriptor se le encargaba describir la posición del Objetivo de manera que el Seguidor pudiera mover su representación del objeto a la misma posición en su pantalla. Luego de una negociación entre ambos jugadores para decidir la mejor posición del objeto, se les asignó a los jugadores una puntuación entre 1 y 100 puntos de acuerdo a qué tan acertado fue el posicionamiento del objetivo por parte del Seguidor.

\begin{figure}
\centering
\includegraphics[scale=0.5]{images/columbia_games.jpg}
\caption{Juego de objetos del Columbia Games}
\label{objects_game}
\end{figure}


Cada sesión consistió de 14 tareas como ésta, cambiando los objetos y sus ubicaciones. En las primeras cuatro tareas, uno de los sujetos tomó el papel del Descriptor; en los siguientes cuatro invirtieron roles, y en las finales seis fueron alternando los roles de Descriptor y Seguidor.

\subsection{Anotaciones sobre comportamiento social}

Varios aspectos del comportamiento de los jugadores durante los juegos de objetos fueron anotados mediante la herramienta de crowdsourcing \emph{Amazon Mechanical Turk}\footnote{https://www.mturk.com}. Cada anotador escuchó el audio correspondiente a una tarea del juego y tuvo que responder a varias preguntas sobre cada uno de los sujetos, entre las que se encuentran:

\begin{itemize}
  \item ¿el sujeto contribuye para el éxito del equipo?  (\svcontributes)
  \item ¿el sujeto parece comprometido con el juego? (\svengaged)
  \item ¿el sujeto se expresa correctamente? (\svclear)
  \item ¿el sujeto piensa lo que va a decir? (\svplanning)
  \item ¿el sujeto alienta a su compañero? (\svencourages)
  \item ¿el sujeto le hace difícil hablar a su compañero? (\svdifficult)
  \item ¿el sujeto está aburrido con el juego? (\svbored)
  \item ¿al sujeto no le agrada su compañero? (\svdislikes)
\end{itemize}

\noindent Cada uno de estos audios fue puntuado por cinco anotadores, que respondieron por sí o por no para cada una de las preguntas. El puntaje que recibe cada una de las preguntas (a las cuales llamaremos a partir de ahora \emph{variables sociales}) consiste en la cantidad de respuestas afirmativas que recibió, teniendo un rango de 0 a 5. Por ejemplo, una tarea dada podría tener puntaje 3 para la variable social `el sujeto A se expresa correctamente' o puntaje 5 para la variable `el sujeto B dirige la conversación'.

\subsection{Extracción de variables acústico/prosódicas}

La herramienta \emph{Praat} \footnote{http://www.fon.hum.uva.nl/praat/} fue utilizada para extraer automáticamente las variables \ap del corpus. Las variables que medimos fueron el tono, la intensidad, la proporción de vocalizaciones, jitter, shimmer, cantidad de sílabas, cantidad de fonemas, y la proporción de ruido sobre armónicos. Estos atributos fueron medidos en cada uno de los segmentos de habla del corpus.

Repasemos algunos conceptos que necesitamos para definir las variables acústicas.

\begin{itemize}
  \item \emph{f0} refiere a la frecuencia fundamental de una onda, que es el recíproco del período de ésta. El \emph{tono} o \emph{pitch} es la percepción que tenemos de la frecuencia fundamental, que nos marca cuán agudos o graves son los sonidos.
  \item \emph{Intensity} refiere al volumen o intensidad de la onda. Ésta mide la amplitud de la onda, y es la percepción de cuán fuerte es el sonido.
  \item \emph{jitter y shimmer} se refieren, en un intervalo de tiempo, a los desplazamientos de la onda de la verdadera periodicidad y de la amplitud, respectivamente. Están asociadas con la percepción de la calidad de la voz.
  \item Un \emph{fonema} es la articulación simple de sonidos del habla, tanto de vocales como de consonantes. Ejemplos de fonemas son los sonidos de las letras u, a, s, k en español.
  \item El \emph{noise-to-harmonics ratio} (abreviado NHR) puede considerarse como una medida de calidad de la voz, que cuantifica la proporción de ruido que hay en ésta.
\end{itemize}


En la siguiente tabla resumimos estas features. Recordemos nuevamente que estas features son medidas en un intervalo de tiempo.

\begin{figure}[h!]
\centering

\begin{tabular} {|c|c|}
  \hline
  Variable & Descripción \\
  \hline
  \hline
  \FOMEAN & Valor medio de la frecuencia fundamental \\\hline
  \FOMAX  & Valor máximo de la frecuencia fundamental \\\hline
  \ENGMEAN & Valor medio de la intensidad \\\hline
  \ENGMAX & Valor máximo de la intensidad \\\hline
  \NOISETOHARMONICS & Noise-to-harmonics ratio \\\hline
  \LOCALSHIMMER & Shimmer medido \\\hline
  \LOCALJITTER  & Jitter medido \\\hline
  \SYLAVG & Cantidad de sílabas por segundo \\\hline
  \PHONAVG & Cantidad de fonemas por segundo \\\hline
\end{tabular}
\end{figure}

\section{Selección de Ventana}
\label{sec:window_selection}
En \cite{KOU2008.2} se menciona una elección de \emph{frame step} y \emph{frame length} de 10s y 20s respectivamente. En el caso de nuestro corpus, quisimos buscar los parámetros que mejor se ajustaban a éste, manteniendo la superposición del 50\% entre ventanas sucesivas. Con lo que nos queda que $FL = 2 * FS$

¿Qué queremos optimizar? La métrica que elegimos para ésto es encontrar un balance entre un frame no tan grande (para no suavizar en exceso la curva) y que nos reduzca considerablemente la cantidad de indefiniciones; es decir, aquellas ventanas que tomamos en un interlocutor que no tienen ninguna interacción de su parte. Para ver ésto, graficamos la cantidad de indefiniciones en función del step tomado.

\begin{figure}
\centering
\includegraphics[width=10cm]{images/window_selection.png}
\end{figure}



Dentro del rango de $FS \in \{5'',6'', \ldots ,15'' \}$, graficamos para cada sesión, tarea y cada interlocutor las curvas de indefiniciones. A su vez, para mayor claridad, graficamos una curva que promedie todas las tareas de una sesión.


Para tener una visión general de lo que ocurría en todas las sesiones, graficamos una curva promedio de todas las sesiones. En ésta puede observarse que hasta $8''-10''$ hay un fuerte descenso de las indefiniciones, que luego se atenúa. Dado que en general tenemos tareas cortas, preferimos tomar $8''$ como step, y $16''$ como largo de ventana.

OBS: podríamos cambiar ésto a un boxplot!

\begin{figure}
\centering
\includegraphics[height=5cm]{images/window_selection_for_session.png}
\end{figure}

\section{Time plots}
\label{sec:time_plots}
Usando la técnica descripta con las variaciones que consideramos en la anterior sección, generamos dos series de tiempo para cada tarea. Como antes mencionamos, la ventana elegida es de 16 segundos con un step de 8 segundos lo cual da un overlap del 50\%.

Dada una ventana, puede ocurrir que alguno de los interlocutores no haya hablado, o su interacción haya sido demasiado breve como para medir sus variables \ap. Como ya mencionamos en la sección \ref{sec:tama_modifications}, y a diferencia de \cite{KOU2008.2}, construimos las series sin ese punto, y sin interpolarlo tampoco.

De estas tareas, sólo nos quedamos con aquellas que tengan al menos 5 puntos definidos para cada serie, de manera que tenga sentido poder calcular la correlación cruzada más adelante. Con esto, no sólo nos interesa la duración de la charla, sino cierta calidad de las series generadas. En la tabla \ref{tab:time_series} pueden verse las tareas que tuvimos en consideración, junto a su duración.

\begin{table}
\centering

\adjustbox{max width=\textwidth}{\begin{tabular}{|l|llllllllllll|}
\toprule
Task & S-01 & S-02 & S-03 & S-04 & S-05 & S-06 & S-07 & S-08 & S-09 & S-10 & S-11 & S-12 \\
\hline
01 &         -- &         -- &    149.8   &         -- &         -- &         -- &         -- &         -- &     54.5   &    106.0   &         -- &     56.1   \\
02 &         -- &         -- &         -- &         -- &         -- &         -- &         -- &         -- &     41.7   &     63.8   &         -- &         -- \\
03 &         -- &     51.7   &         -- &     80.7   &     77.9   &     69.2   &     68.4   &     49.6   &         -- &    122.2   &     81.0   &         -- \\
04 &         -- &    187.2   &     93.3   &     76.1   &     79.9   &     99.2   &     84.3   &         -- &     58.0   &    129.6   &     67.9   &     95.2   \\
05 &         -- &         -- &         -- &     86.3   &         -- &    126.7   &    145.8   &     90.7   &     45.7   &    134.2   &         -- &         -- \\
06 &         -- &         -- &         -- &         -- &         -- &    148.2   &     50.6   &     60.2   &     46.1   &     66.7   &     46.7   &     40.2   \\
07 &         -- &     66.0   &         -- &    117.7   &         -- &     72.4   &         -- &     87.7   &     85.9   &    110.6   &     65.7   &         -- \\
08 &         -- &    458.8   &     98.6   &    203.8   &         -- &    188.7   &     59.9   &     48.1   &         -- &    157.4   &         -- &     81.1   \\
09 &         -- &         -- &         -- &     75.5   &    134.2   &     83.0   &    108.7   &         -- &     62.1   &    404.0   &     41.0   &     92.5   \\
10 &     50.1   &    231.3   &    162.8   &    242.5   &         -- &    122.4   &     71.1   &     74.7   &         -- &    356.0   &     69.8   &     92.7   \\
11 &         -- &     74.4   &         -- &     98.6   &     70.1   &         -- &     58.9   &         -- &     72.9   &    104.0   &     59.4   &    101.9   \\
12 &     61.3   &     90.1   &    129.1   &    182.9   &         -- &    130.3   &     75.8   &     57.6   &         -- &    101.6   &         -- &     64.8   \\
13 &     55.1   &    124.0   &    108.1   &    144.1   &    114.7   &         -- &         -- &     83.8   &     94.0   &    174.0   &     84.8   &     91.5   \\
14 &         -- &     75.3   &         -- &         -- &    107.3   &         -- &     52.5   &    144.3   &     75.5   &    108.4   &     91.6   &     98.4   \\
\bottomrule
\end{tabular}
}

\caption{Tabla de tareas seleccionadas y sus duraciones}
\label{tab:time_series}
\end{table}


Como primer paso siempre recomendado en el análisis de series de tiempo \cite{CHATFIELD}, graficamos los time plots conjunto de cada par de series, a la vez que sus autocorrelogramas (ver apéndice \ref{sec:time_series}). En la figura \ref{fig:time_plot} podemos observar un ejemplo de esto.

A priori, las series tienen aspecto de series autoregresivas de orden uno. Es decir, series que son de la forma $X_t = \alpha X_{t-1} + e_t + c$, con $e_t$ ruido blanco, $\alpha$ y $c$ constantes. Este hecho es esperable  por la construcción misma del método TAMA, ya que la ventana de cada punto tiene un solapamiento con la ventana anterior. Más aún, uno esperaría que $\alpha \sim 0.5$ ya que nuestras ventanas tienen ese índice de overlap. Los autocorrelogramas de las series, por otro lado, tienen en su mayoría un valor significativo en $k = 1$, el valor del $\alpha$ de la autoregresión.


\begin{figure}
\centering
\includegraphics[width=15cm]{images/time_plot_with_autocorrelation.png}
\caption{Time-plot generado por el método TAMA, junto a su autocorrelograma}
\label{fig:time_plot}
\end{figure}

El hecho de que los autocorrelogramas desciendan rápidamente a cero es un indicio de que las series de tiempo construidas son estacionarias. Esto nos habilita a efectuar el análisis bivariado de las series.

%\subsection{Corpus}
\begin{frame}
  \frametitle{Columbia Games Corpus}
  \framesubtitle{Descripción}
  \begin{columns}
    \column{0.35\textwidth}
      \begin{figure}
        \includegraphics[width=\textwidth]{images/columbia_games_color.jpg}
      \end{figure}

    \column{0.65\textwidth}

    \begin{itemize}
      \item Desarrollado por Agustín Gravano para su tesis doctoral.
      \item Corpus de conversaciones diádicas en Inglés Norteamericano
      \item 12 sesiones con 14 tareas/juegos cada una.
      \item En cada sesión, se sentó a dos participantes en una cabina profesional de grabación, y una cortina opaca colgando entre ellos para evitar la comunicación visual.
      \item Los participantes contaron con computadoras a través de las cuales interactuaban mediante juegos.
    \end{itemize}
  \end{columns}
\end{frame}


\begin{frame}
  \frametitle{Columbia Games Corpus}
  \framesubtitle{Anotaciones sociales}

  Cinco anotadores escucharon el audio correspondiente a una tarea del juego y respondieron a varias preguntas sobre los sujetos:


\begin{table}
\adjustbox{max width=0.8\textwidth}{
\begin{tabular}{|l|l|}
  \hline

  \textbf{Nombre} & \textbf{Pregunta} \\
  \hline

  \svcontributes &  ¿el sujeto contribuye para el éxito del equipo?  \\ \hline
  \svengaged &  ¿el sujeto parece comprometido con el juego? \\ \hline
  \svclear &  ¿el sujeto se expresa correctamente? \\ \hline
  \svplanning &  ¿el sujeto piensa lo que va a decir? \\ \hline
  \svencourages &  ¿el sujeto alienta a su compañero? \\ \hline
  \svdifficult &  ¿el sujeto le hace difícil hablar a su compañero? \\ \hline
  \svbored &  ¿el sujeto está aburrido con el juego? \\ \hline
  \svdislikes &  ¿al sujeto no le agrada su compañero? \\ \hline
\end{tabular}
}
\end{table}

De cada una de estas preguntas obtenemos un puntaje de 0 a 5, para cada hablante de cada tarea.
\end{frame}


\begin{frame}
\frametitle{Extracción de features acústico-prosódicas}

Usando el software Praat \footnote{http://www.fon.hum.uva.nl/praat/} se extrajeron las variables acústico-prosódicas para cada segmento de habla

\begin{table}
\adjustbox{max width=0.8\textwidth}{
\centering
\begin{tabular} {|c|c|}
  \hline
  Variable & Descripción \\
  \hline
  \FOMEAN & Valor medio de la frecuencia fundamental \\\hline
  \FOMAX  & Valor máximo de la frecuencia fundamental \\\hline
  \ENGMEAN & Valor medio de la intensidad \\\hline
  \ENGMAX & Valor máximo de la intensidad \\\hline
  \NOISETOHARMONICS & Noise-to-harmonics ratio \\\hline
  \LOCALSHIMMER & Shimmer medido \\\hline
  \LOCALJITTER  & Jitter medido \\\hline
  \SYLAVG & Cantidad de sílabas por segundo \\\hline
  \PHONAVG & Cantidad de fonemas por segundo \\\hline
\end{tabular}
}
\end{table}
\end{frame}



\subsection{Modificaciones a TAMA}

\begin{frame}
\frametitle{TAMA}
\framesubtitle{Nuestras modificaciones}

\begin{itemize}
  \item Usamos un step de 8s y un tamaño de ventana de 16s, manteniendo el solapamiento del 50\%.
  \item A diferencia del trabajo original de Kousidis, utilizamos series con datos faltantes.
  \item Pero sólo nos quedamos con aquellas que tengan 5 o más puntos definidos.
\end{itemize}
\end{frame}


\begin{frame}
\frametitle{TAMA}
\framesubtitle{Nuestra medida de mimetización}

  \begin{columns}
    \column{0.33\textwidth}
    \begin{figure}[t]
      \includegraphics[scale=0.28]{images/time_plot.png}
    \end{figure}

    \begin{figure}[t]
      \includegraphics[scale=0.28]{images/cross_correlogram_2.png}
    \end{figure}

    \column{0.66\textwidth}

    Definimos dos medidas de mimetización

    \begin{align*}
      \fwentrainment{AB}^{(1)} &= r_s \text{ con s maximizando } |r_k|,  k <= 0  \\
      \fwentrainment{BA}^{(1)} &= r_s \text{ con s maximizando } |r_k|,  k >= 0  \\
      \fwentrainment{XY}^{(2)} &= |\fwentrainment{XY}^{(1)}|
    \end{align*}

    Segunda métrica motivada por estudios sobre la antimimetización. Healey et al (2014) sugiere que puede ser una conducta de adaptación cooperativa.

    Levitan et al (2015) da más indicios en esa dirección.
  \end{columns}
\end{frame}

\subsection{Análisis de la relación con variables sociales}

\begin{frame}
\frametitle{Mimetización y relación con variables sociales}

  Para analizar la relación entre las variables sociales ($V$) y nuestras medidas de \emph{mimetización} ($\mathcal{E}$), planteamos un modelo de regresión lineal.

    \begin{equation}
      V_i \sim \beta_1 + \beta_2 \mathcal{E}_i
    \end{equation}

  Nuestra hipótesis es

  \begin{enumerate}
    \item Si $V$ es una variable de carácter positivo, entonces $\beta_2 > 0$
    \item Si $V$ es una variable de carácter negativo, entonces $\beta_2 < 0$
  \end{enumerate}
\end{frame}

\section{Panel de datos}
\label{sec:panel_data}

Luego de construir las series de tiempo para cada una de las conversaciones que seleccionamos anteriormente, pasamos a construir una gran tabla que se utilizó en los experimentos de regresión. Para condensar todos nuestros datos, armamos una tabla por cada variable a/p, que contiene información definida para cada interlocutor, tarea y sesión de nuestro corpus.

Cada columna de esta tabla representa los datos de un interlocutor dentro de una tarea. Éste hecho lo usamos fuertemente a la hora de definir los grupos en nuestro modelo de Efectos Fijos. En la figura \ref{fig:panel_data}

\begin{figure}
\centering
\adjustbox{max width=\textwidth}{
\begin{tabular}{|c|c|}
  \hline
  \emph{Campo} & \emph{Descripción} \\\hline
  session & número de sesión \\\hline
  speaker & 0 si corresponde al interlocutor A; B en otro caso \\\hline
  task & número de tarea \\\hline
  count & La cantidad de puntos definidos que tiene la serie \\\hline
  entrainment & Si $speaker=0$, es $A\rightarrow B$; $B \rightarrow A$ en otro caso \\\hline
  best\_lag & el lag del cross-correlogram donde se logra el \emph{entrainment} \\\hline

  engaged\_in\_game & ¿el sujeto parece comprometido con el juego? \\\hline
  difficult\_for\_partner\_to\_speak & ¿el sujeto dirige la conversación? \\\hline
  contributes\_to\_successful\_completion & ¿el sujeto contribuye para el éxito del equipo? \\\hline
  gives\_encouragement & ¿el sujeto alienta a su compañero?\\\hline
  making\_self\_clear & ¿el sujeto se expresa correctamente?\\\hline

  planning\_what\_to\_say & \\\hline
  bored\_with\_game & ¿el sujeto se muestra aburrido? \\\hline
  dislikes\_partner &  ¿al sujeto no le agrada su compañero? \\\hline
\end{tabular}}
\label{fig:panel_data}
\end{figure}



esto nos da una tabla de 210 filas, y 21 columnas. Para cada sesión y speaker, podemos pensar que tenemos una serie de tiempo donde el tiempo es cada tarea, y los datos son el entrainment y las variables sociales. En la jerga econométrica, llamamos a este tipo de datos \emph{de panel}\cite{gujarati1999}: un conjunto de mediciones temporales sobre un mismo sujeto a lo largo del tiempo. En este caso el sujeto es un interlocutor en una sesión, el tiempo son las tareas, y las mediciones son los entrainments


\begin{figure}
\centering
\adjustbox{max width=\textwidth}{
\begin{tabular}{lrrrrrrrrr}
\toprule
session &  speaker &  task &  entrainment &  bored\_with\_game &  engaged\_in\_game &  gives\_encouragement &  making\_self\_clear &  planning\_what\_to\_say \\
\midrule
1 &        0 &    10 &     0.581475 &                0 &                5 &                    5 &                  5\\
1 &        0 &    12 &    -0.569677 &                1 &                5 &                    5 &                  5\\
1 &        0 &    13 &     0.533701 &                2 &                4 &                    5 &                  4\\
1 &        1 &    10 &    -0.917101 &                0 &                5 &                    2 &                  3\\
1 &        1 &    12 &     0.467112 &                0 &                5 &                    4 &                  2\\
1 &        1 &    13 &    -0.602364 &                0 &                5 &                    4 &                  3\\
2 &        0 &     3 &     0.520696 &                0 &                4 &                    5 &                  5\\
2 &        0 &     4 &    -0.241060 &                0 &                5 &                    4 &                  4\\
2 &        0 &     7 &     0.743719 &                0 &                5 &                    4 &                  5\\
2 &        0 &     8 &     0.147362 &                0 &                5 &                    4 &                  2\\
\bottomrule
\end{tabular}
}


\caption{Panel de datos}
\label{panel_data}
\end{figure}


\begin{figure}
\centering
\includegraphics[width=10cm]{images/linear_regression.jpg}
\caption{Ejemplo de Regresión Lineal}
\end{figure}

\section{Análisis de regresión}

Llegado a este punto, dada una variable a/p, nos interesaría evaluar la relación entre el entrainment y las distintas variables sociales. Con esto en mente, planteamos un modelo de regresión lineal donde nuestra variable explicativa será la mimetización, y la variable \emph{dependiente} será la variable social.

En base a ésto, podremos observar cuál es la variación conjunta de ellas. Es esperable que, al aumentar la mimetización, aumenten ciertas variables sociales (por ejemplo, la compenetración en el juego) y que otras desciendan (el aburrimiento).


\subsection{Modelo clásico de Regresión Lineal}

En el modelo clásico de regresión lineal, tenemos un conjunto de valores fijos $X_1, X_2, \ldots, X_n$, que son llamadas variables independientes. Asociado a cada uno de estos valores fijos, tenemos variables aleatorias $Y_1, \ldots, Y_n$. Asumimos, además, que nuestras variables son de la forma

\begin{equation}
  Y_i = E[Y|X_i] + u_i
\end{equation}

donde $u_i$ es la perturbación estocástica de la variable.

Asumiendo que $E[Y|X_i]$ es una función lineal de $X_i$; es decir, que existen $\beta_1, \beta_2 \in \mathbb{R}$ que cumplen

\begin{equation}
  E[Y|X_i] = \beta_1 + \beta_2 X_i
\end{equation}

obtenemos que

\begin{equation}
  Y_i = \beta_1 + \beta_2 X_i + u_i
\end{equation}

Nuestro objetivo es poder entonces conseguir estimadores $\widehat{\beta_1}, \widehat{\beta_2}$ que nos permitan analizar y predecir el comportamiento conjunto de estas variables.

\subsection{Nuestro modelo}

Sea entonces una variable acústica/prosódica (por ejemplo, el pitch o la intensidad), y una variable social de las que acabamos de enumerar en \ref{sec:panel_data}. Sean $E_1, \ldots, E_n$ los valores de entrainment para el set de datos que definimos en \ref{sec:panel_data}, y sean $V_1, V_2, \ldots V_n$ los valores de la variable social de cada conversación.

Sobre éstas variables es que planteamos nuestro modelo de regresión lineal clásica: queremos ver qué relación hay tomando como variable ``fija'' al entrainment, y como variable dependiente a la variable social. Queremos hallar, entonces $\widehat{\beta_1}, \widehat{\beta_2} \in \mathbb{R}$

\begin{equation}
  V_i \simeq \widehat{\beta_1} + \widehat{\beta_2} E_i
\end{equation}


Para ello, calcularemos los estimadores $\widehat{\beta_1}, \widehat{\beta_2} \in \mathbb{R}$ mediante el método \emph{QR} (insertar referencia aquí) que nos provee el lenguaje R. A su vez, luego de ésto efectuaremos un análisis de significancia sobre $\beta_2$ para verificar que sean distintos de 0.


Uno esperaría que un alto \emph{entrainment} se relacione con un alto valor de ciertas variables sociales \cite{BRE1996}, por ejemplo la compenetración con el juego, el ayudar a terminarlo. Esto significa esperar que el valor de $\widehat{\beta_2}$; y se relacione con bajos valores de otras, como el aburrimiento, o el rechazo percibido hacia el compañero.


\subsection{Modelo agrupado o \emph{pooled}}

En el modelo agrupado o \emph{pooled}, no distinguimos entre datos provenientes de distintos ``grupos'' \cite{gujarati1999} y sobre éstos calculamos la regresión lineal, agrupando todos los datos disponibles.

Un problema que surge con este tipo de regresión es que niega todo tipo de \emph{heterogeneidad} de los datos: estos pueden provenir de interlocutores más o menos empáticos, o cuya interacción en el juego se vio influída por factores no medidos en el experimento. Todo ésto es descartado, aún cuando puede afectar seriamente  el resultado obtenido.

AGREGAR GRAFICO DE EJEMPLO PARA ESTO

\subsection{Modelo de Efectos Fijos dentro de cada grupo}


El modelo de efectos fijos agrega el concepto de heterogeneidad permitiendo que cada sujeto tenga su propio valor de ordenada al origen. En términos formales, reformulemos nuestro modelo de la sección anterior:

\begin{equation}
  V_{it} = \beta_1 + \beta_2 * E_{it} + u_{it}
\end{equation}

donde $i$ es la cantidad de sujetos (en nuestro caso, 12 sesiones por 2 interlocutores = 24), $t$ es la variable de tiempo (en nuestro caso, las tareas de cada sesión). El modelo de efectos fijos nos permite


DEFINIR SUJETO



\chapter{Regresión Lineal Agrupada}
En esta sección, mostraremos el primer experimento realizado. Este consistió en aplicar un modelo de regresión lineal de cada variable social sobre el \entrainment, sin desagregar los datos por sesión y hablante.

Una variación que usamos en el presente experimento (y en el posterior) es utilizar como variable dependiente el valor absoluto del \entrainment, en base a estudios que sugieren que los interlocutores pueden también diferenciarse como un rasgo positivo en la conversación.

\section{Nuestro modelo de regresión}

Dada una variable acústica/prosódica (por ejemplo, el pitch o la intensidad), queremos investigar la relación entre entrainment y las distintas variables sociales medidas. Sea $V$ una variable social de las enumeradas en la tabla \ref{sec:panel_data}. Sean $E_1, \ldots, E_n$ los valores de entrainment para el set de datos que definimos en la sección \ref{sec:panel_data}, y sean $V_1, V_2, \ldots V_n$ los valores de la variable social de cada conversación.

Sobre éstas variables es que planteamos nuestro modelo de regresión lineal clásica, para analizar qué relación hay tomando como variable ``fija'' al \entrainment, y como variable dependiente a la variable social. El problema, entonces, es hallar estimadores $\widehat{\beta_1}, \widehat{\beta_2} \in \mathbb{R}$ de modo que 

\begin{equation}
  V_i \simeq \estintercept + \estslope E_i
  \label{eq:my_model}
\end{equation}


Para ello, calculamos los estimadores mediante el método \emph{QR} que nos provee el lenguaje R. A su vez, luego de esto efectuamos un análisis de significancia sobre $\estslope$ para verificar que sea estadísticamente distintos de 0.

Uno esperaría que un alto \emph{entrainment} se relacione con un alto valor de ciertas variables sociales, por ejemplo la compenetración con el juego, el ayudar a terminarlo. En términos de la ecuación \ref{eq:my_model}, esperamos que $\estslope \geq 0$. De manera inversa, cuando las variables sociales tienen un carácter negativo de la conversación, esperamos que $\estslope \leq 0$. 


El modelo de regresión que usamos en este primer experimento se denomina agrupado o \emph{pooled}, ya no distinguimos entre datos provenientes de distintos ``grupos'' \cite{gujarati1999} y sobre estos calculamos la regresión lineal, agrupando todos los datos disponibles.

Un problema que surge con este tipo de regresión es que niega todo tipo de \emph{heterogeneidad} de los datos: estos pueden provenir de interlocutores más o menos empáticos, o cuya interacción en el juego se vio influída por factores no medidos en el experimento. Todo esto es descartado, aún cuando puede afectar seriamente  el resultado obtenido.

En el siguiente capítulo ahondaremos un poco más en cómo definimos los grupos en nuestro trabajo.

\section{Resultados sobre \entrainment}

Los resultados de este experimento no fueron interesantes ya que dieron valores de $\estslope$ con muy baja signifcancia. En \ref{fig:regresion_clasica} puede verse el gráfico de regresión lineal del \entrainment contra distintas variables sociales, tomando como variable acústico-prosódica a \FOMEAN. En las figuras \ref{fig:regresion_clasica_tabla_1} y \ref{fig:regresion_clasica_tabla_2} podemos observar la tabla de las regresiones, con las estimaciones obtenidas y sus valores de significancia para todas las variables sociales. Se observa que no sólo las pendientes tienen un muy bajo valor absoluto, sino que además ni siquiera tienen los signos que esperábamos en un principio.

En base a lo arrojado por este análisis de regresión, intentamos introducir variaciones del experimento. La primera, es cambiar la variable explicativa por el valor absoluto del entrainment.


\begin{figure}[b!]
\includegraphics[width=15cm]{images/regression_F0_MEAN_1.pdf}
\caption{Gráfico de los pares entrainment-variable a/p, junto a la regresión lineal obtenida para \emph{F0\_MEAN}}
\label{fig:regresion_clasica}
\end{figure}


% "ENG_MEAN"
% latex table generated in R 3.2.2 by xtable 1.8-0 package
% Thu Jan  7 03:01:56 2016
\begin{tabular}{rrrrr}
  \hline
 & Estimate & Std. Error & t value & Pr($>$$|$t$|$) \\
  \hline
bored\_with\_game & 10.65 & -0.44 & 0.00 & 0.66 \\
  difficult\_for\_partner\_to\_speak & 10.56 & -0.54 & 0.00 & 0.59 \\
  contributes\_to\_successful\_completion & 59.62 & -1.15 & 0.00 & 0.25 \\
  engaged\_in\_game & 73.14 & 0.55 & 0.00 & 0.59 \\
  gives\_encouragement & 47.49 & -0.04 & 0.00 & 0.97 \\
  making\_self\_clear & 52.97 & -0.99 & 0.00 & 0.33 \\
  planning\_what\_to\_say & 32.02 & -1.81 & 0.00 & 0.07 \\
  dislikes\_partner & 9.61 & -0.94 & 0.00 & 0.35 \\
   \hline
\end{tabular}

% "F0_MEAN"
% latex table generated in R 3.2.2 by xtable 1.8-0 package
% Thu Jan  7 03:01:56 2016
\begin{tabular}{rrrrr}
  \hline
 & Estimate & Std. Error & t value & Pr($>$$|$t$|$) \\
  \hline
bored\_with\_game & 10.63 & -0.18 & 0.00 & 0.86 \\
  difficult\_for\_partner\_to\_speak & 10.68 & -0.97 & 0.00 & 0.33 \\
  contributes\_to\_successful\_completion & 59.38 & 0.89 & 0.00 & 0.37 \\
  engaged\_in\_game & 73.41 & 0.77 & 0.00 & 0.44 \\
  gives\_encouragement & 47.59 & 1.09 & 0.00 & 0.28 \\
  making\_self\_clear & 52.91 & -0.32 & 0.00 & 0.75 \\
  planning\_what\_to\_say & 31.49 & 0.39 & 0.00 & 0.70 \\
  dislikes\_partner & 9.81 & -1.74 & 0.00 & 0.08 \\
   \hline
\end{tabular}

% "F0_MAX"
% latex table generated in R 3.2.2 by xtable 1.8-0 package
% Thu Jan  7 03:01:56 2016
\begin{tabular}{rrrrr}
  \hline
 & Estimate & Std. Error & t value & Pr($>$$|$t$|$) \\
  \hline
bored\_with\_game & 11.08 & 2.46 & 0.00 & 0.01 \\
  difficult\_for\_partner\_to\_speak & 10.65 & 0.52 & 0.00 & 0.60 \\
  contributes\_to\_successful\_completion & 60.08 & -0.98 & 0.00 & 0.33 \\
  engaged\_in\_game & 74.20 & -0.57 & 0.00 & 0.57 \\
  gives\_encouragement & 48.17 & -1.04 & 0.00 & 0.30 \\
  making\_self\_clear & 53.86 & -2.32 & 0.00 & 0.02 \\
  planning\_what\_to\_say & 31.86 & -1.05 & 0.00 & 0.30 \\
  dislikes\_partner & 9.65 & 0.97 & 0.00 & 0.33 \\
   \hline
\end{tabular}

% "NOISE_TO_HARMONICS_RATIO"
% latex table generated in R 3.2.2 by xtable 1.8-0 package
% Thu Jan  7 03:01:56 2016
\begin{tabular}{rrrrr}
  \hline
 & Estimate & Std. Error & t value & Pr($>$$|$t$|$) \\
  \hline
bored\_with\_game & 10.75 & -0.44 & 0.00 & 0.66 \\
  difficult\_for\_partner\_to\_speak & 10.64 & 0.16 & 0.00 & 0.87 \\
  contributes\_to\_successful\_completion & 60.34 & -0.75 & 0.00 & 0.45 \\
  engaged\_in\_game & 74.49 & -0.16 & 0.00 & 0.87 \\
  gives\_encouragement & 48.39 & -0.79 & 0.00 & 0.43 \\
  making\_self\_clear & 53.60 & -0.11 & 0.00 & 0.91 \\
  planning\_what\_to\_say & 31.99 & -0.02 & 0.00 & 0.99 \\
  dislikes\_partner & 9.62 & -1.39 & 0.00 & 0.17 \\
   \hline
\end{tabular}





\nota{Mover esto a antecedentes}
\section{Valor absoluto de \entrainment}

En nuestra definición de \entrainment en el contexto de series de tiempo, la definimos como el valor de la correlación cruzada (en un sentido de los lags) con mayor valor absoluto. Esto puede dar, como resultado, valores positivos entre 0 y 1 a los cuales consideramos como \entrainment; o bien valores negativos entre -1 y 0, estos considerados como anti-\entrainment: la divergencia de las features a/p medidas a través del tiempo.

Este fenómeno de anti-\entrainment o antimimicry \cite{CHAR1999} refiere al proceso por el cual uno de los interlocutores no imita al otro sino más bien todo lo contrario, acentúa alguna diferencia. Si bien estudios de larga data como \cite{bourhis1973language} o \cite{dabbs1969similarity} lo emparentan con una connotación negativa, \cite{healey2014divergence} y \cite{levitan2015acoustic} sugieren que puede entenderse este fenómeno como una conducta de adaptación cooperativa. No sólo éso, sino que este fenómeno de mimetización complementaria es más prevalente que la mimetización a secas \cite{levitan2015acoustic}.

En base esto es que decidimos probar alguna medida que capture positivamente el fenómeno de la anti-mimetización de igual manera que con el \entrainment antes definido. Es decir, esperamos que cuando tengamos o bien \entrainment o \entrainment complementario ocurra que tenemos valores altos de variables sociales de carácter positivo. Mutatis mutandis con las variables sociales de connotación negativa.

Con este fin, en vez de utilizar sólo el valor de \entrainment como variable explicativa, efectuamos el mismo análisis pero utilizando el valor absoluto del \entrainment como tal. Usar esto permite captar y valorar el \entrainment complementario de la misma manera que el ``positivo'' y valorar su relación con las variables sociales medidas.

\section{Resultados sobre \absentrainment}

\nota{Escribir acá, y poner tablas sobre \absentrainment en pooled}



\chapter{Regresión Lineal con Efectos Fijos}
\section{Modelo de Efectos Fijos}

\newcommand{\slopeestim}[1] { $\estslope \sim #1$ }

El modelo de efectos fijos agrega el concepto de heterogeneidad permitiendo que cada sujeto tenga su propio valor de ordenada al origen. En el caso concreto de nuestro corpus, dicha heterogeneidad puede deberse a multiplicidad de factores no medidos en él. Por ejemplo, la personalidad de los sujetos es un factor no medido y que puede influír en la dinámica de la interacción entre éstos.

Como el entrainment que medimos mediante el proceso TAMA es un proceso direccional (es decir, medimos tanto la influencia de un interlocutor sobre el otro y viceversa), definimos los ``grupos'' como las observaciones de entrainment y sus variables sociales de una sesión y el interlocutor sobre el cual consideramos la direccionalidad, tal y cual lo definimos en una sección anterior. Ésto nos arroja la cantidad de 24 grupos de observaciones



\section{Definición Formal de Efectos Fijos}

NOTA: Escribir ésto

\section{Resultados}

El modelo de efectos fijos sobre el valor absoluto del \emph{entrainment} dio valores sustancialmente apreciables. Las variables a/p \ENGMAX, \FOMEAN y \NOISETOHARMONICS poseen valores altamente significativos ( p-valor menor a 0.05) para la regresión con efectos fijos para al menos 2 variables sociales. En la tabla \ref{regresion_efectos_fijos_tabla} podemos ver la tabla del test de coeficientes con las variables sociales significativas resaltadas. En \ref{sign_table} podemos observar una tabla de doble entrada que simplifica la interpretación de las tablas de coeficientes, marcando aquellos pares de variables a/p y variables sociales con coeficientes significativos y su signo.




\begin{tabular}{rrrrr}
  \hline
 \ENGMAX & Estimate & Std. Error & t value & Pr($>$$|$t$|$) \\
  \hline
contributes\_to\_successful\_completion & 0.0497 & 0.4262 & 0.1165 & 0.9074 \\
  \myhighlight making\_self\_clear & 1.6581 & 0.3864 & 4.2909 & 0.0001 \\
  engaged\_in\_game & 0.3307 & 0.2576 & 1.2840 & 0.2008 \\
  planning\_what\_to\_say & 0.5005 & 0.5327 & 0.9395 & 0.3487 \\
  gives\_encouragement & 0.4264 & 0.3792 & 1.1246 & 0.2622 \\
  \myhighlight difficult\_for\_partner\_to\_speak & -0.7200 & 0.2858 & -2.5190 & 0.0126 \\
  bored\_with\_game & 0.2163 & 0.2560 & 0.8450 & 0.3992 \\
  dislikes\_partner & -0.4318 & 0.3443 & -1.2541 & 0.2114 \\
   \hline
\end{tabular}

\begin{tabular}{rrrrr}
  \hline
 \FOMEAN & Estimate & Std. Error & t value & Pr($>$$|$t$|$) \\
  \hline
 \myhighlight contributes\_to\_successful\_completion & 1.0274 & 0.3025 & 3.3962 & 0.0008 \\
  \myhighlight making\_self\_clear & 0.8307 & 0.3934 & 2.1115 & 0.0361 \\
  \myhighlight engaged\_in\_game & 0.8850 & 0.2750 & 3.2182 & 0.0015 \\
  planning\_what\_to\_say & 0.7167 & 0.5400 & 1.3273 & 0.1860 \\
  gives\_encouragement & 0.0075 & 0.3941 & 0.0190 & 0.9848 \\
  difficult\_for\_partner\_to\_speak & -0.5975 & 0.3928 & -1.5209 & 0.1300 \\
  \myhighlight bored\_with\_game & -0.7586 & 0.2481 & -3.0572 & 0.0026 \\
  dislikes\_partner & 0.0371 & 0.3800 & 0.0977 & 0.9223 \\
   \hline
\end{tabular}

\begin{tabular}{rrrrr}
  \hline
 \NOISETOHARMONICS & Estimate & Std. Error & t value & Pr($>$$|$t$|$) \\
  \hline
\myhighlight contributes\_to\_successful\_completion & 0.7041 & 0.3404 & 2.0686 & 0.0400 \\
  \myhighlight making\_self\_clear & 1.3344 & 0.3537 & 3.7725 & 0.0002 \\
  engaged\_in\_game & 0.0954 & 0.3462 & 0.2756 & 0.7832 \\
  planning\_what\_to\_say & -0.1874 & 0.4177 & -0.4485 & 0.6543 \\
  gives\_encouragement & 0.7234 & 0.4782 & 1.5127 & 0.1321 \\
  difficult\_for\_partner\_to\_speak & -0.1941 & 0.3436 & -0.5648 & 0.5729 \\
  \myhighlight bored\_with\_game & 0.5876 & 0.3028 & 1.9404 & 0.0539 \\
  dislikes\_partner & 0.3582 & 0.3330 & 1.0755 & 0.2835 \\
   \hline
\end{tabular}

\begin{figure}[ht]
\centering
% psl is "Positive Slope"
\newcommand{\psl} { $+$ }
\newcommand{\ppsl} { $++$ }
\newcommand{\pppsl} { $+++$ }

% nsl stands for "Negative SLope"
\newcommand{\nsl} { $-$ }
\newcommand{\nnsl} { $--$ }
\newcommand{\nsl} { $---$ }


\begin{tabular}{| c | c | c | c | c | c |}
  \hline
               &\ENGMAX  & \ENGMEAN  & \FOMEAN  & \FOMAX  & NOISERATIO  \\
  \hline
  contributes  &         &           & \psl     &         & \psl        \\ \hline
  clear        & \pppsl  &           & \psl     &         & \psl        \\ \hline
  engaged      &         &           & \psl     &         &             \\ \hline
  planning     &         &           &          &         &             \\ \hline
  encourages   &         &           &          &         &             \\ \hline
  difficult    & \nsl    &           &          &         &             \\ \hline
  bored        &         &           & \nsl     &         &             \\ \hline
  dislikes     &         &           &          &         &             \\ \hline
  \hline
& PHON\_AVG & PHON\_COUNT & SHIMMER & SYL\_AVG & SYL\_COUNT \\
  \hline
contributes  &      &  &  &  &        \\ \hline
  clear      & \psl &  &  &  & \psl   \\ \hline
  engaged    &      &  &  &  &        \\ \hline
  planning   &      &  &  &  &        \\ \hline
  encourages &      &  &  &  &        \\ \hline
  difficult  &      &  &  &  &        \\ \hline
  bored      &      &  &  &  &        \\ \hline
  dislikes   &      &  &  &  &        \\ \hline
  \hline
\end{tabular}


\caption{Tabla que representa los resultados significantes del experimento. En una de las entradas, tenemos los nombres abreviados de las variables sociales, y en la otra las variables a/p. El símbolo \psl representa valor significante y positivo de la pendiente de la regresión de efectos fijos, mientras que \nsl representa significante y negativo }

\label{sign_table}

\end{figure}



\chapter{Conclusiones y trabajo futuro}

%%%% BIBLIOGRAFIA
\chapter{Bibliografía}
\backmatter
\bibliographystyle{alpha}
\bibliography{tesis}
\end{document}
